\begin{definition}{Online Bipartite Matching with Recourse}
	The Recourse model is similar to the Semi-Online model in structure, but different in function. One anticlique $U$ is still always known, and the other, $V$, still arrives in an online fashion, but the goal of the algorithm is to maintain a maximal matching with as few changes to $M$ as possible. The sum of these changes over all arriving vertices is the algorithm's Recourse, a way of measuring the algorithm like competitiveness. The goal when creating an algorithm is to minimize Recourse
\end{definition}

we will begin the introduction with a simple proof
\begin{remark}
	In the Online Edge Arrival Recourse model, where individual edges arriven in an online fashion and all vertices are visible, no algorithm has a recourse below $O\Omega(n^2)$
\end{remark}
\begin{proof}
	The graph is a simple even path $(v_1,\dots,v_{2k})$. First edge $(v_k,v_k+1)$ arrives and are matched. Then, for $i\in(1,\dots,k-1)$, edges $(v_{k-i},v_{k-i+1}),(v_{k+i},v_{k+i+1})$ arrive sequentially, effectively extending the path from the front and back. Once every two edge arrivals the maximal matching changes and becomes the set of all edges except the previous matching, forcing the algorithm to change the entire matching. Thus, a total of $\Omega(n^2)$ changes is required
\end{proof}
so the distinction between the arrival of individual edges and the arrival of vertices, and all edges connected to them, is very important, since
\begin{theorem}
	There is an algorithm, SAP, that in has a Recourse of $O(n \log^2 n)$
\end{theorem}
we will also showcase the proof for a similar lower bound on vertex arrival
\begin{theorem}
	No algorithm has a Recourse below $\Omega(n\log n)$ in the Recourse model
\end{theorem}

\section{Lower Bound}
\input{sections/recourse/lowerbound.tex}

\section{SAP Algorithm}
the algorithm presented, as well as the proof follow the ideas from \cite{obmRecourse}
\begin{algorithm}
	\caption{SAP}
	\label{alg:sap}
	\begin{algorithmic}
		\item on Arrival of vertex $v$: Augment the matching with the shortest augmenting path from $v$ is exists
	\end{algorithmic}
\end{algorithm}

The crucial part of this proof is the proof of the following
\begin{lemma}[Augmentation Lemma]
	\label{recourse-crucial}
	$\forall h\in\mathbb{Z}$, SAP, over it's entire runtime, augments the matching with a total of no more then $\frac{4n\ln(n)}{h}$ augmenting paths longer than $h$
\end{lemma}
First, we will prove SAP's correctness

\begin{lemma}
	SAP maintains a maximum matching
\end{lemma}
\begin{proof}
	We will conduct the proof by induction over time $t$. Time at $t=0$ when no vertices have arrived forms a trivial base case with the maximum matching being $M=\emptyset$. For the induction step:
	Knowing
	\begin{enumerate}
		\item At time $t$, the matching $M_t$ is maximum, and
		\item A matching is maximum iff there are no augmenting paths
	\end{enumerate}
	we conclude that if $M_t$ isn't maximum at time $t+1$ then there must be an augmenting path that didn't exist at time $t$, meaning that the newly arriving $v$ must be it's starting point
\end{proof}

Then, we'll prove the Recourse bound using \autoref{recourse-crucial}
\begin{proof}
	We will sum all possible $h$ augmenting path lengths by grouping them into intervals $l_i=[2^i,2^{i+1})$\\
	Paths in $l_i$ contain no more than $2^{i+1}$ edges and there are no more than $\frac{4n\log n}{2^i}$ in each group, giving $O(n\log n)$ edges per group.
	Since there are $\log_2(n)+1$ groups, total Recourse, which is no greater than the sum of all edges in all paths in all groups is no greater than $O(n \log^2n)$
\end{proof}

and finally, we will present the proof of \autoref{recourse-crucial} in the following sections. We will first define a Server Flow and prove several properties in the first section to prove \autoref{recourse-crucial} in the second section

\subsection{Server Flow}
In this subsection we will prove the Augmentation Lemma

A server flow is a fractional assignment from $V$ (clients) to $U$ (servers), in which an online vertex sends flow, up to 1, to offline vertices (servers), every one of which has a unique capacity to receive flow.
\begin{definition}[Server Flow]
	A server flow is $\alpha$ is a map from $U$ to $\mathbb{R}_{\geq 0}$ such that there are non-negative reals $\{x_e|e\in E\}$ such that
	\begin{align*}
		\forall v\in V:\sum\limits_{u\in N(v)}x_{uv}=1 && \forall u\in U:\sum\limits_{v\in N(u)}x_{uv}=\alpha(u)
	\end{align*}
\end{definition}
we refer to these $x$-values as realizing the flow. Note that if $\forall u:\alpha(u)=1$ then this is a fractional matching.\\
Going forward we assume $\forall v\in V:\sum\limits_{u\in N(v)}x_{uv}\geq 1$ as otherwise no server flow is possible

Informally, a Balanced Server Flow is one where every client sends flow to the neighboring servers that receives the least flow
\begin{definition}[Balanced Server Flow]
	A Server Flow is Balanced if
	\begin{align*}
		A(v)= \underset{u\in N(v)}{\mathrm{argmin}}\alpha(u)\\
		\forall v\in V\forall u\in N(v)\backslash A(v):x_{uv}=0
	\end{align*}
\end{definition}

where $A(v)$ are called active neighbors of $v$

\begin{lemma}
	\label{balanced-sf-iff}
	A balanced server flow exists iff $\forall v\in V: |N(v)|\geq 1$
\end{lemma}
\begin{proof}
	Proving requirement (if) simple, if an online vertex doesn't have a neighbor, it's resource cannot be received, and there is no server flow.\\
	We will prove uniqueness by using convex programming. The program
	\begin{align*}
		\text{minimize} & \sum\limits_{u\in U}\alpha(u)^2\\
		\text{subject to} & \sum\limits_{u\in N(v)}x_{uv}=1 & \forall v\in V\\
		& \sum\limits_{v\in N(u)}x_{uv}=\alpha(u) & \forall u\in U
	\end{align*}
	is convex, and moreover, the objective function is strictly convex (due to being a $n$-dimensional paraboloid) meaning that the program has a unique (with regard to $\alpha$-values) optimum\\
	Observe that the optimum is a balanced server flow: constraint ensure that it is a server flow and if it wasn't balanced there would be an online vertex $v$ that sends non-zero flow to $u,u'$ and $\alpha(u)<\alpha(u')$. Notice that the objective function can be decreased by increasing $\alpha(u)$ and reducing $\alpha(u')$ until they are identical, therefore the flow must be balanced.
\end{proof}

We will now prove that the balanced server flow is unique (regarding $x$-values) using an auxiliary lemma

%insert above, unify proofs
\begin{lemma}
	\label{balanced-sf-aux}
	For a balanced server flow $x_{uv}$
	Let $\alpha_u=\sum\limits_{v\in V}x_{uv}$ be server flow of $u$\\
	Let $x'_{uv}$ be feasible server flow\\
	Let $\alpha'_u=\sum\limits_{v\in V}x'_{uv}$ be resulting server load\\
	\begin{align*}
		\sum\limits_{u\in U}\alpha_u^2\leq\sum\limits_{u\in U}\alpha_u\alpha'_u
	\end{align*}
\end{lemma}
\begin{proof}
	for $v\in V$ let $\mu(v)=\underset{u\in N(v)}{\min}\alpha_u$ be the minimum server load of $u\in N(v)$
	\begin{align*}
		    &\sum\limits_{u\in U}\alpha_u^2
		   =&\sum\limits_{u\in U}\sum\limits_{v\in V}x_{uv}\alpha_u
		   =&\sum\limits_{v\in V}\sum\limits_{u\in U}x_{uv}\alpha_u
		   =&\sum\limits_{v\in V}\sum\limits_{u\in U}x_{uv}\mu(v)\\
		   =&\sum\limits_{u\in U}x_{uv}\mu(v)
		   =&\sum\limits_{v\in V}\sum\limits_{u\in U}x'_{uv}\mu(v)
		\leq&\sum\limits_{v\in V}\sum\limits_{u\in U}x'_{uv}\alpha_u
		   =&\sum\limits_{u\in U}\sum\limits_{v\in V}x'_{uv}\alpha_u\\
		   =&\sum\limits_{u\in U}\alpha'_u\alpha_u
	\end{align*}
	with the inequality being caused by the fact that the flow $x_{uv}$ is balanced, so $\alpha_u\geq\mu(v)$ for any edge $(u,v)\in E$
\end{proof}

\begin{lemma}
	\label{balanced-sf-unique}
	The server flow in \autoref{balanced-sf-iff} if it exists
\end{lemma}
Do note that we do not claim that the corresponding $x$-values are unique, only the corresponding $\alpha$
\begin{proof}
	Using \autoref{balanced-sf-aux} we can argue that the balanced server flow is the optimum of the linear program in the proof of \autoref{balanced-sf-iff}.\\
	We can treat a feasible solution $\alpha'$ as a vector, with \autoref{balanced-sf-aux} showing that $||\alpha||\leq \alpha\alpha'$ which with non-negativity implies $||\alpha||\leq||\alpha'||$\\
	Therefore, any balanced server flow is the optimum of the convex program, which in turn implies that it must be unique since the optimum is unique.
\end{proof}

Let $\alpha$ be the balanced server flow at time $t$ and $\alpha'$ be the balanced server flow at time $t+1$, after a new online vertex arrives.\\
$\Delta\alpha(u)=\alpha(u)'-\alpha(u)$
\begin{lemma}
	\label{balanced-sf-insert}
	$\forall u\in U:\Delta\alpha(s)\geq 0$
\end{lemma}
\begin{proof}
	Let $U'=\set{u|\Delta\alpha(u)<0}$\\
	We will prove by contradiction that $U'=\emptyset$\\
	Assume $U'\neq\emptyset$, let $\alpha^\star=\underset{u\in U'}{\alpha'(u)}$.\\
	Let $U^-=\set{u\in U|\alpha(u)\leq\alpha^\star}$\\
	Let $U^\Delta=\set{u\in U|\alpha(u)>\alpha'\land\alpha'(u)=\alpha^\star}$\\
	Let $U^+=\set{u\in U|\alpha(u)>\alpha'\land\alpha'(u)>\alpha^\star}$\\
	These three partition $U$, $\emptyset\neq U^\Delta\subseteq U'$\\
	Let $V^\Delta$ be the set of all $v\in V$ with an active neighbor at time $t$. Since each sends one unit of flow, $\sum\limits_{u\in U}\alpha(u)\leq|V^\Delta|$\\
	The server flow being balanced at $t$ implies there is no edges from $V^\Delta$ to $S^-$ at time $t$.\\
	Further, there are no active edges from $V^\Delta$ to $S^+$ at time $t+1$\\
	Therefore, all active edges incident to $V^\Delta$ must connect to $U^\Delta$ implying $\sum\limits_{u\in U}\alpha'(u)\geq|V^\Delta|$.
	This in turn, when combined with the previous inequality implies $\sum\limits_{u\in U}\alpha'(u)\geq \sum\limits_{u\in U}\alpha(u)$\\
	This is a contradiction as, by definition of $U^\Delta$, $\sum\limits_{u\in U}\alpha'(u)< \sum\limits_{u\in U}\alpha(u)$ 
\end{proof}

\begin{lemma}
	\label{balanced-sf-affects}
	when $v$ arrives, $\forall u\in U:\alpha(u)<\underset{u'\in N(v)}{\min}\alpha(u')\implies\Delta\alpha(u)=0$
\end{lemma}
Or, more informally ans simplistically, arrival of $v$ only connected to $s$ cannot affect the flow of offline vertices whose flow is less than the flow among neighbors of $s$
\begin{proof}
	consider the balanced server flow at time $t$ (before inserting $v$)\\
	Let $U^+=\set{u\in U|\alpha(u)\geq\underset{u'\in N(v)}{\min}\alpha(u')}$\\
	Ler $U^-=U\backslash U^+$ be complement of $U^+$\\
	Let $V^+=\set{v\in V|N(v)\subseteq U^+}$. Note that, before the insertion of $v$, there are no edges between $V^+$ and $U^-$, together with the fact that there are no active edges from $V^-$ to $U^+$ because the flow is balanced implying $\sum\limits_{u\in U}\alpha(u)=|V^-|$\\
	After the insertion of $v$, by definition of $U^-$ $v$ has no neighbors in $U^-$, maintaining previous property, implying the new balanced flow $\alpha'$ has the property $\sum\limits_{u\in U}\alpha'(u)\leq |V^-|$, meaning that $\sum\limits_{u\in U}\Delta\alpha(u)\leq 0$\\
	If, for some $u_1$, $\Delta\alpha(u_1)>0$ then $\Delta\alpha(u_2)<0$ for some other $u_2$, contradicting \autoref{balanced-sf-insert}
\end{proof}

\begin{lemma}
	\begin{align*}
		\forall T\subseteq U: |\set{v\in V|A(v)\subseteq T}|\leq \sum\limits_{u\in T}\alpha(u)\leq |\set{v\in V|A(v)\cap T\neq\emptyset}|
	\end{align*}
\end{lemma}
\begin{proof}
	The first inequality is a direct result of the fact that every online vertex in the first set contributes exactly $1$ to the sum\\
	The second inequality is a direct result of the fact that every online vertex contributes exactly $1$ to $\sum\limits_{u\in U}\alpha(u)$, and in the inequality every online vertex contributing anything to $\sum\limits_{u\in T}\alpha(u)$ is counted as if contributing $1$
\end{proof}

\begin{lemma}
	Let $\overline{\alpha}=\underset{u\in U}{\max}\alpha(u)$ be maximal necessity, $\overline{U}=\set{u\in U|\alpha(u)=\overline{\alpha}}$ maximally necessary servers, $K'=\set{v\in V|A(v)\cap\overline{S}\neq\emptyset}$ active neighbors, then $N(K')=\overline{U}$, $|K'|=\overline{\alpha}|\overline{U}|$
\end{lemma}
\begin{proof}
	Let $K=\set{v\in V|A(v)\subseteq\overline{U}}$\\
	Note that $K\subseteq K'$\\
	Note that, by definition of $\overline{U}$, $K'\neq\emptyset$ and $v\in K'\implies N(v)=A(v)\subseteq\overline{U}$\\
	Therefore $K=K'$, $N(K')=\overline{U}$.\\
	By the previous lemma $|K'|=|K|\leq\overline{\alpha}|\overline{U}|\leq|K'|$ implying equality
\end{proof}

\begin{lemma}
	\label{balanced-sf-max}
	\begin{align*}
		\overline{\alpha}=\underset{\emptyset\subset K\subseteq V}{\max}\frac{|K|}{|N(K)|}
	\end{align*}
	and $\forall K\subseteq V:|K|=\overline{\alpha}|N(K)|\implies\forall u\in N(K)\alpha(u)=\overline{\alpha}$
\end{lemma}
\begin{proof}
	Definition of server flow implies, for $K\subseteq V$ $|K|\leq\sum\limits_{u\in N(K)}\alpha(s)\leq\overline{\alpha}|N(K)|$ implying $|K|\leq\overline{\alpha}|N(K)|$.\\
	Let $K'$ be defined as in above lemma. Then $\overline{\alpha}=\frac{|K'|}{N(K')}\leq\underset{\emptyset\subset K\subseteq V}{\min}\frac{|K|}{|N(K)|}\leq\overline{\alpha}$\\
	If $|K|=\sum\limits_{u\in N(K)}\alpha(u)=\overline{\alpha}|N(K)|$ then $\forall u\in U:\alpha(u)\leq\overline{\alpha}$ implies $\forall u\in N(K):\alpha(u)=\overline{\alpha}$
\end{proof}

\subsection{Replacement Analysis}
We will now draw a connection between Server Flow and Matching and use the previous lemmas to prove \autoref{recourse-crucial}

\begin{lemma}
	$G=(U,V,E)$ contains a matching of size $|V|$ iff $\forall u\in U\alpha(u)\leq 1$
\end{lemma}
\begin{proof}
	The lemma follows directly from \autoref{balanced-sf-max}: $\forall K\subseteq V:|K|\leq |N(K)|$ iff $\overline{\alpha}\leq 1$.
\end{proof}

Define $V_M\subseteq V$ as follows: when $v$ arrives, $V'$ are the vertices that arrived previously. $v\in V_M$ iff the maximum matching after $v$ arrived is greater than the one before.\\
$G_M=(U,V_M,E)$ as opposed to the final graph $G=(U,V,E)$\\
$\alpha_M$ is the balanced server flow in $G_M$. By \autoref{balanced-sf-unique} it is unique
An augmenting tail from vertex $v\in U\cup V$ is an alternating path starting at $v$ and ending at an unmatched $u\in U$. Tail is active if all edges that aren't in the matching are active

\begin{lemma}[Expansion]
	\label{recourse-expansion}
	Assume $\alpha_M(u)=1-\epsilon$ for some $u\in U,\epsilon>0$. There is an active augmenting tail for $u$ of length at most $\frac{2}{e}\ln(|V_M|)$
\end{lemma}
\begin{proof}
	Notice that any offline vertex $u'$ reachable from $u$ by an active augmenting tail has $\alpha_M(u')\leq 1-\epsilon$\\
	Let $K_i$ be online vertices $v$ such that there is an augmenting tail $u_0,v_0,\dots,v_{k-1},u_k$ where $u=u_0,v=v_j$ for some $j<i$. Do note that $v\in K_i\implies v\in K_{i+1}$.
	\begin{align*}
		|K_i|\leq\sum\limits_{u'\in A(K_i)}\alpha_M(u')\leq\sum\limits_{u'\in A(K_i)}(1-\epsilon)=|A(K_i)|(1-\epsilon)
	\end{align*}
	Implying $|A(K_i)|\geq\frac{|K_i|}{1-\epsilon}$.
	Assuming no active augmenting tails from $u$ of length at least $2(i-1)$, every offline vertex in $A(K_i)$ is matched and the set of offline vertices they are matched to is $K_{i+1}$\\
	Bijection $A(K_i)\leftrightarrow K_{i+1}$ is implied by the perfect matching, so $|K_{i+1}|=|A(K_i)|$ therefore $|V_M|\geq|K_{i+1}|\geq\left(\frac{1}{1-\epsilon}\right)^i|K_{1}|=\left(\frac{1}{1-\epsilon}\right)^i$. This, alongside $1-e\leq e^{-epsilon}$ imply that $i\leq\frac{\ln|V_M|}{\ln\frac{1}{1-\epsilon}}\leq\frac{1}{\epsilon}\ln|V_M|$.
	Therefore, for any $i>\frac{1}{\epsilon}\ln|V_M|$ there is an active augmenting tail of length at most $2(i-1)$ implying the lemma
\end{proof}

We can now proceed to the proof of \autoref{recourse-crucial}. Recall the lemma statement: By, on vertex arrival, adding augmenting paths from the arriving vertex if possible, no more than $\frac{4n\ln n}{h}$ augmenting paths of length greater than $h$ are added in total

\begin{proof}
	$n=|V|\geq|V_M|$. For $h\leq 4\ln n$ this lemma holds trivially since there are at most $n$ augmenting paths in total. We therefore assume $h>4\ln n$ for the remainder of the proof.\\
	Note that any augmenting path must be in $G_M$\\
	Let $V'\subseteq V_M$ be the online vertices whose shortest augmenting paths are longer than $h+1$ when they are added, the lemma is equivalent to the statement $|V'|\leq \frac{4n\ln n}{h}$\\
	For $v\in V'$ the shortest augmenting tail from each $u\in N(v)$ has length at least $h$, so by \autoref{recourse-expansion} each has $\alpha_M(u)\geq 1-\frac{2\ln n}{h}$.\\
	Let $U'$ be all online vertices $u$ that some point have $\alpha_M(u)\geq1-\frac{2\ln n}{h}$. By \autoref{balanced-sf-insert} $U'$ is also the set of offline vertices $u$ that $\alpha_M(u)\geq1-\frac{2\ln n}{h}$ in the final graph\\
	By \autoref{balanced-sf-affects} $v\in V'$ implies arrival of $v$ only increases the flow of offline vertices in $U'$ that had a flow not below $1-\frac{2\ln n}{h}$ before arrival.\\
	Since $G_M$ has a perfect matching, $\forall u\in U:\alpha_M(u)\leq\alpha(u)\leq 1$, implying the flow can only increase by no more than $2n\ln n$\\
	Since each $v\in V$ contributes exactly $1$ flow, the total number of these online vertices is $|V'|<|U'|\frac{2\ln n}{h}$\\
	Again, each $v\in V_M$ sends $1$ flow implying $n\geq|V_M|\geq(1-\frac{2\ln n}{h})|U'|>\frac{|U'|}{2}$ since $h>4\ln n$.
	This results in $|V'|<|U'|\frac{2\ln n}{h}<\frac{4n\ln n}{h}$ proving the lemma
\end{proof}

\section{Implications}
Finally, we come to the connection between this model and the more classical Semi-Online Matching, and the idea of strengthening deterministic algorithms. From these results we can clearly extrapolate that in a model where an algorithm can make a certain amount of corrections, a correction count of $O(n\ln^2 n)$ is enough for a deterministic algorithm to find a perfect matching and a correction count less then $O(n\ln n)$ is not enough